\chapter{Introduction}
Imbuing intelligent agents with the capability for autonomous-decision-making in real-world environments is challenging. For decades, reinforcement learning has offered a robust mathematical framework for acquiring optimal behavioral skills through interactions within an environment to maximize cumulative rewards. The integration of deep neural networks as high-capacity function approximators has propelled RL even farther, leading to remarkable achievements across diverse domains, from mastering complex games to enabling basic robotic locomotion and manipulation.
However, the widespread success of modern machine learning, as seen in areas such as computer vision and natural language processing, fundamentally hinges on the availability of large and highly diverse datasets. As seen in Figure \ref{fig:online_offline}, in conventional, online RL, this data is collected through a continuous interaction loop: the agent executes its current policy in the environment, collects the resulting experiences into a buffer, and then uses this fresh data to update its policy. This cycle of exploration and optimization is repeated until the policy converges. With this paradigm, the acquisition of such extensive data becomes prohibitively impractical for real-world applications. For example, autonomous driving would need to practice in numerous cities for every model update, which is costly, time-consuming, let alone being unsafe. This inherent limitation created a gulf of generalization between lab-bound RL successes and the complex, open-world settings where intelligent decision-making is most critically needed.

\begin{figure}[H]
  \centering
  \includegraphics[width=0.45\textwidth]{online_rl}
  \includegraphics[width=0.45\textwidth]{offline_rl}
  \caption{Comparison of Online and Offline Reinforcement Learning.}
  \label{fig:online_offline}
\end{figure}

Offline Reinforcement Learning emerges as a pivotal paradigm to bridge this gap. Offline RL focuses on enabling the training of policies solely from previously collected datasets, without additional online interaction. This approach holds immense potential for transforming vast amounts of historical data into robust, generalizable decision-making engines, much like how supervised learning revolutionized pattern recognition. The implications are profound across various fields from robotics, healthcare, and autonomous driving to logistics, finance, and operations research, as it offers a viable pathway to deploy sophisticated AI systems in domains where active data collection is unfeasible, costly, or unsafe.

Finally, the ability to leverage large and diverse datasets in an offline setting aligns the field of robotics learning with the broader paradigm of foundation models, as seen in computer vision and natural language processing. This opens the door to pre-training generalist agents on massive, heterogeneous datasets to acquire a broad repertoire of skills, which can then be efficiently fine-tuned for specific downstream tasks. This paradigm thus offers a compelling vision for the future: transforming vast archives of passive, observational data into active, embodied intelligence.

\section{Objectives}

The primary aim of this dissertation is to develop and validate a novel, deterministic framework for imitation learning in offline reinforcement learning, leveraging the principles of flow matching. The specific objectives established to achieve this aim are as follows:
\begin{enumerate}
	\item To design and validate a systematic methodology for curating offline reinforcement learning datasets from actively trained agents.
	\item To design and implement "Flow Planner," a modular framework that utilizes flow matching for the generation of trajectories from offline data.
	\item To analyze the impact of various conditioning mechanisms on the precision and generalizability of the trajectories generated by the framework.
	\item To evaluate the performance of the "Flow Planner" framework, benchmarking its core attributes such as inference speed and policy effectiveness against diffusion models.
\end{enumerate}

\section{Achievements}
The primary contribution of this dissertation is the development and validation of Flow Planner, a novel hierarchical framework for imitation learning, which is composed of two independent technical achievements:

\begin{enumerate}
	\item A Novel Formulation of Flow-Based Trajectory Planning: This work introduces the first specialization of the Conditional Flow Matching (CFM) framework for the domain of robotic trajectory planning. By reformulating the CFM objective to generate different sequences of states and actions, this contribution establishes a new class of deterministic, high-speed generative planners that directly address the latency limitations of contemporary diffusion-based methods.
	\item A Hierarchical Framework for Long-Horizon Control: The second key contribution is the Flow Planner architecture, a new two-stage hierarchical framework that decouples long-horizon strategic reasoning from low-level reactive control. This design solves the scaling and coherence problems inherent in monolithic, end-to-end planners by factorizing the planning problem into a strategic Waypoint Planner and a reactive Action Planner, a novel application for flow-based models.
\end{enumerate}

These core contributions are supported by:

\begin{enumerate}
	\item A Comprehensive Analysis of Flow-Based Planners: This work provides an in-depth empirical analysis of applying flow-matching models to robotics, offering insights into the role of different model architectures and conditioning mechanisms, alongside a direct comparison to the diffusion paradigm.
	\item A Reproducible Data Curation Pipeline and Benchmark Datasets: This work contributes a flexible and extensible pipeline for curating high-quality offline reinforcement learning datasets. This pipeline was used to generate three novel benchmark datasets, which have been publicly released alongside their corresponding trained expert agents to facilitate reproducibility and future research.	
\end{enumerate}
These contributions culminate in a robust and efficient framework for imitation learning, with the findings currently being prepared for submission to the IEEE International Conference on Intelligent Robots and Systems (IROS).

\section{Overview of Dissertation}

The remainder of this dissertation is structured as follows:
\begin{itemize}
	\item Chapter 2 provides the theoretical foundation for this work, first examining established methods for addressing distribution shift problems, then exploring contemporary solutions that leverage generative models.
	\item Chapter 3 presents the core technical foundation of this dissertation. It details various approaches to dataset modeling with comprehensive mathematical formulations. The chapter ends by introducing "Flow Planner," the proposed hierarchical planning framework, and provides implementation specifics using Meta's flow matching library.
	\item Chapter 4 details the methodology for data curation and introduces the experimental environments. This includes a technical discussion on the implementation of the data pipeline using Gymnasium wrappers and the Minari library for dataset management.
	\item Chapter 5 provides a comprehensive analysis and discussion of the empirical results. This includes an investigation of various conditioning mechanisms and a benchmark comparison of the framework's performance against baseline models.
	\item Chapter 6 concludes the dissertation by summarizing the key contributions, discussing the implications of the findings, and outlining several promising directions for future research.
\end{itemize}





















